\hypertarget{chapter-8-robust-command-design}{%
\chapter{Robust Command Design}\label{chapter-8-robust-command-design}}

The preceding chapters covered how macros are defined, how expansion behaves, and how to give a macro a clean calling interface. This chapter is about a different concern: making sure a macro keeps behaving correctly once it leaves the context it was first tested in, gets used by someone other than its author, or has to coexist with dozens of other macros in a project that has been growing for years. None of the individual techniques here are difficult; the discipline is in applying them consistently across an entire project rather than only where a bug has already shown up.

\hypertarget{fragile-commands-still}{%
\section{Fragile Commands, Still}\label{fragile-commands-still}}

Chapter 2 introduced fragile commands as one of the places LaTeX's abstraction visibly leaks: a macro that works correctly the first time it is used but breaks when its argument is written to an auxiliary file and reprocessed, because that reprocessing happens in a context where full expansion is required and the macro was not built to survive it. This might seem like a legacy concern from an older, pre-LaTeX3 era, but it has not gone away in modern engines. Section titles, captions, and running headers are still written to auxiliary files and reread; a macro used inside any of these that is not built from fully expandable material, all the way down through anything it calls internally, will still fail in exactly the way it failed decades ago, regardless of whether the surrounding document uses LuaLaTeX, XeLaTeX, or the classic pdfTeX engine.

The practical response has two parts. The first is \texttt{\textbackslash{}protect}, which prevents a fragile command from being expanded prematurely in a moving-argument context, deferring its expansion to the point where it is actually being typeset rather than the point where it is being recorded into an auxiliary file. The second, and generally the better long-term fix where it is available, is writing the macro to be robust in the first place, built entirely from expandable primitives or from \texttt{\textbackslash{}newcommand}, \texttt{\textbackslash{}NewDocumentCommand}, and other LaTeX-kernel definition commands, which produce robust macros by default in current LaTeX releases. \texttt{\textbackslash{}protect} remains necessary for older or lower-level macro definitions and for third-party commands whose internals are not under the author's control; for anything written from scratch for a new project, preferring robust-by-construction definition commands avoids the problem before it starts.

\hypertarget{namespacing}{%
\section{Namespacing}\label{namespacing}}

A large project accumulates macros quickly: helper commands for notation, shortcuts for repeated phrases, formatting wrappers around theorem environments, and macros supporting the book's own custom document class. Left unnamespaced, this accumulation eventually collides, either with a macro name the author reused without noticing, or with a name a loaded package happens to define for its own purposes, silently overriding one definition with the other and producing a failure that surfaces far from its actual cause.

The LaTeX3 convention of \texttt{\textbackslash{}ClassName\_module\_function:nn}-style naming, using underscores and colons reserved for package and class code rather than document-level macros, is one solution, but it is heavier than most single-author book projects need throughout. A lighter convention that captures most of the benefit is prefixing project-specific macros with a short, consistent tag tied to the project or the book's own notation system: \texttt{\textbackslash{}rsvpfield}, \texttt{\textbackslash{}rsvpboundary}, rather than \texttt{\textbackslash{}field} and \texttt{\textbackslash{}boundary}, which are far more likely to collide with something a loaded package already defines. This costs a small amount of typing and pays for itself the first time a package update introduces a name that would otherwise have silently shadowed a document-level macro with the same name, a failure mode that is often diagnosed by symptom rather than by cause unless the naming convention is disciplined enough to make the collision itself unlikely.

For a project spanning multiple related books sharing a common notation system, this namespacing matters even more, since macros defined for one book's specific needs are liable to be copied into a companion project's preamble, and an unprefixed macro name is far more likely to already mean something different in the second project than a prefixed one is.

\hypertarget{writing-for-a-future-reader}{%
\section{Writing for a Future Reader}\label{writing-for-a-future-reader}}

The macros in a large project are read far more often than they are written: read while debugging, read while extending, read by a collaborator, read by the same author eighteen months later with no memory of the reasoning behind a particular implementation choice. A macro's definition should be legible on its own terms, which in practice means a small set of habits applied consistently. A comment above any macro whose purpose is not obvious from its name and arguments, stating what it is for rather than restating what the code already shows syntactically. Meaningful argument names in \texttt{xparse} and \texttt{l3keys} definitions where the tooling allows it, rather than relying on positional \texttt{\#1}, \texttt{\#2} with no accompanying explanation of what each position means. Consistent formatting of macro definitions across a project, so that a reader's eye can find the argument list, the body, and any embedded conditionals in the same relative position every time, rather than having to re-parse an unfamiliar layout for every new macro encountered.

None of this is about elegance for its own sake. A project that will be actively maintained, extended, and rebuilt for years benefits from exactly the same discipline that a long-lived software codebase benefits from, and for the same reason: the cost of an illegible macro is paid not once, at the moment it is written, but repeatedly, every time it has to be understood again by someone who is not currently holding its full context in mind.

Part III closes here, with the tools for writing macros that behave predictably and interfaces that are pleasant to call. Part IV turns to typography and fonts: the layer where LuaLaTeX and XeLaTeX's modern capabilities are put to direct use, after three chapters spent entirely on how a document is built and programmed rather than on how it looks.
